\subsection{THE SET OF NATURAL INTEGERS}

DEFINITION 1. A set is said to be infinite if it is not finite.

In particular, a cardinal is infinite if it is not an integer.

The relation ``there exists an infinite set'' implies that the relation ``x is an integer'' is collectivizing (Chapter II, § 1, no. 4); for if a is an infinite cardinal and n an arbitrary integer, we cannot have \(a \leqslant n\) (§ 4, no. 2, Proposition 2). We have therefore n \textless{} a for all integers n, which shows that the set of integers \textless{} a (§ 3, no. 2, Remark following Theorem 1) contains all integers. Conversely, if the relation ``x is an integer'' is collectivizing, then the set E of integers is an infinite set. For each integer n, the interval {[}0, n{]} is a subset of \(n + 1\) elements of E (§ 5, no. 3, Proposition 5). Therefore Card (E) \(\geqslant n + 1 > n\) . But to say that \(\text{Card}(E) \neq n\) for every integer n means that E is infinite.

We now introduce the following axiom :

A5 (``Axiom of infinity''.) There exists an infinite set.

It is not known whether or not this axiom can be deduced from the axioms and axiom schemes introduced previously; although the question has not been definitively settled, it is to be presumed that this axiom is independent of the others.

The preceding remarks then prove the following theorem :

THEOREM 1. The relation ``x is an integer'' is collectivizing.

We shall denote by N the set of integers (also called ``the set of natural integers'' when necessary to avoid ambiguity). The cardinal of N is denoted by \(\aleph_0\) . Whenever \(\mathbf{N}\) is considered as an ordered set, it is always the ordering (called the usual ordering) defined in §3, no. 2 that is under consideration, unless the contrary is expressly stated.

DEFINITION 2. A sequence (resp. a sequence of elements of a set E) is a family (resp. a family of elements of E) whose index set is a subset of N. The sequence is said to be infinite if its index set is an infinite subset of N.

Let \(P\{n\}\) be a relation and let I denote the set of integers n such that \(P\{n\}\) is true. I is then a subset of N. A sequence \((x_n)_{n \in \mathbf{I}}\) is then sometimes written \((x_n)_{P\{n\}}\) , and \(x_n\) is called the nth term in the sequence. A sequence whose index set is the set of integers \(n \geqslant k\) is often written \((x_n)_{k \leqslant n}\) or \((x_n)_{n \geqslant k}\) , or even just \((x_n)\) if k = 0 or k = 1. Under the same

conditions, for example, the notations \(\prod_{\mathbf{P}\nmid n}\mathbf{X}_n\) and \(\prod_{n = k}\mathbf{X}_n\) are used to denote the product of a sequence of sets \((\mathbf{X}_n)_{n\in \mathbf{I}}\) , and there are analogous notations for unions, intersections, cardinal products, and cardinal sums.

Every subfamily of a sequence is a sequence, called a subsequence of the given sequence.

Two sequences \((x_{n})_{n\in I}, (y_{n})_{n\in I}\) with the same index set are said to differ only in the order of their terms if there exists a permutation \(f\) of the index set \(I\) such that \(x_{f(n)} = y_{n}\) for all \(n\in I\) .

A multiple sequence is a family whose index set is a subset of a product \(N^{p}\) (p an integer) (``double sequence'' for p = 2, ``triple sequence'' for p = 3, and so on).

Let I be a set equipotent to N and let f be a bijection of N onto I. For each family \((x_{t})_{t\in I}\) indexed by the set I, the sequence \(n\to x_{f(n)}\) is said to be obtained by arranging the family \((x_{t})_{t\in I}\) in the order defined by f.~The sequences which correspond in this way to two distinct bijections of N onto I differ only in the order of their terms. For a finite family indexed by a set I of n elements we may similarly define a finite sequence with \([1,n]\) or \([0,n-1]\) as index set, by arranging the family in the order defined by a bijection of one or the other of these intervals onto I.

